\section{Introduction}
\label{sec:intro}


Quantum computing offers a computational model distinct from classical computing by operating on qubits rather than binary bits. Each qubit is represented as a quantum state with probability amplitudes (e.g., $|\psi\rangle = \alpha|0\rangle + \beta|1\rangle$), where superposition enables multiple states, and entanglement creates correlations between qubits. These properties enable parallelism beyond classical computation~\cite{renner2022computational, miguel2023enhancing}. Despite these advantages, current quantum computers remain limited to Noisy Intermediate-Scale Quantum (NISQ) devices with high error rates and short coherence times~\cite{lau2022nisq, saki2019study, preskill2018quantum}. To validate algorithms without hardware-induced errors, classical computers are used for quantum circuit simulation. Specifically, full state vector simulation evaluates all probability amplitudes to provide the exact state distributions and precise gradients essential for training quantum machine learning models~\cite{miller2025phase2, kim2023evidence, xu2025herculean}. 

While full state vector simulation guarantees exact results, its exponential memory footprint forces existing frameworks to rely on state-of-the-art infrastructures, including HPC and cloud systems equipped with terabytes of memory and high-performance GPU clusters~\cite{xu2024atlas, 10.1145/3771577, zhang2021hyquas}. However, as quantum applications transition to real-world deployments (e.g., decentralized quantum sensor networks~\cite{guo2020distributed}, privacy-sensitive medical IoT~\cite{barua2024data}, and autonomous robotics~\cite{skolik2022quantum}), relying on such centralized infrastructures becomes impractical for data-intensive edge workloads~\cite{fernandez2019pre, caro2022generalization}. First, regulated data (e.g., HIPAA-protected medical telemetry, GDPR-protected biometrics) is legally restricted from leaving the device~\cite{theodos2020health, khan2021federated}. Second, field-deployed nodes operate beyond reliable network coverage~\cite{wang2020convergence, lim2020federated}. Third, relying on centralized high-end computing resources incurs prohibitive infrastructure costs across thousands of IoT devices~\cite{xu2024iot, gill2019transformative}. Finally, emerging applications such as hybrid quantum-classical machine learning and Variational Quantum Algorithms (VQAs) over IoT sensor streams require thousands of iterative circuit executions~\cite{verma2025quantum, cerezo2021variational, mastroianni2024variational}. Offloading this optimization process to the cloud incurs network round-trips on every iteration, where communication latencies quickly accumulate into hours of idle blocking time. Therefore, quantum parameters must be optimized on the device that generates the data~\cite{zhang2025breaking, li2025infscaler}.


However, edge-native quantum simulation is fundamentally constrained by an exponential memory wall. The footprint of a quantum state vector grows exponentially, rapidly exceeding the strict physical memory limits of edge SoCs (e.g., 8 GB on an NVIDIA Jetson Orin Nano~\cite{jetson_orin_nano}). Furthermore, unlike servers equipped with discrete GPUs, edge devices employ a Unified Memory Architecture (UMA) where the CPU, GPU, and operating system share a single DRAM pool. When a massive state vector exhausts the shared environment, it starves essential system workspaces, leading to execution bottlenecks and out-of-memory (OOM) failures.


%To overcome this memory barrier, existing approaches attempt to deploy high-performance frameworks (e.g., cuQuantum) using out-of-core memory management. However, applying these frameworks to edge environments exposes an architecture mismatch. In contrast to discrete GPUs, edge devices employ a Unified Memory Architecture (UMA) where the CPU and GPU share a constrained physical memory pool. When standard frameworks attempt to allocate large state vectors within this shared pool, they exhaust the available DRAM. Lacking edge-aware memory management, they fall back on operating system virtual memory (OS-native swapping) to spill data to external NVMe storage, which triggers I/O thrashing. Figure~\ref{intro_motivation_b} shows the impact of this mismatch on a Variational Quantum Eigensolver (VQE) workload. CPU-only execution (Cirq) avoids GPU memory exhaustion but is computationally expensive, requiring over 100 seconds (109,168ms) for a single 29-qubit iteration. Conversely, GPU-accelerated execution (cuQuantum) degrades as it triggers OS swapping near physical memory limits (96.6ms at 29 qubits), before crashing with an Out-of-Memory (OOM) error at 30 qubits due to GPU workspace overheads and CUDA allocator constraints. In contrast, \EdgeQuantum leverages a UMA-optimized architecture to bypass OS bottlenecks, maintaining millisecond-level latency (2.0~ms) at 29 qubits and completing the 30-qubit out-of-core simulation in 10.3 seconds.

\begin{table}[t]
\caption{Comparison with prior work across key capabilities (TA: Tiered Architecture, CS: Compression Support, MN: Multinode Support, UO: UMA Optimization, RS: Reported Scale).}
\vspace{-.2cm}

\centering
\scriptsize
\begin{tabularx}{\columnwidth}{p{1.8cm}|c|c|c|c|c|c|c}
\toprule
\textbf{Framework} & \textbf{Target} & \textbf{GPU} & \textbf{TA} & \textbf{CS} & \textbf{MN} & \textbf{UO} & \textbf{RS} \\
\midrule
Google Cirq~\cite{cirq} & General &  &  &  & \checkmark &  & 32 \\
Qiskit Aer~\cite{qiskit_aer} & General & \checkmark &  &  &  &  & 32+ \\
Google Qsim~\cite{qsim} & General & \checkmark &  &  & \checkmark &  & 32+ \\
cuQuantum~\cite{bayraktar2023cuquantum} & General & \checkmark &  &  & \checkmark &  & 32+ \\
PennyLane~\cite{pennylane} & General & \checkmark &  &  & \checkmark &  & 32+ \\
HyQuas~\cite{zhang2021hyquas} & HPC system & \checkmark &  &  &  &  & 36 \\
Atlas~\cite{xu2024atlas} & HPC system & \checkmark & &  & \checkmark &  & 36 \\
ScaleQsim~\cite{10.1145/3771577} & HPC system & \checkmark &  &  & \checkmark &  & 42 \\
SnuQS~\cite{park2022snuqs} & Datacenter & \checkmark & \checkmark &  &  &  & 42 \\
BMQSim~\cite{zhang2025bmqsim} & Datacenter & \checkmark & \checkmark & \checkmark & \checkmark &  & 36 \\
\textbf{EdgeQuantum} & \textbf{Edge/IoT} & \checkmark & \checkmark & \checkmark &  & \checkmark & \textbf{39} \\
\bottomrule
\end{tabularx}
\label{intro_table}
\end{table}

Many previous studies, as summarized in Table~\ref{intro_table}, optimize quantum circuit simulation for different computing infrastructures, ranging from general-purpose workstations to HPC systems. Standard frameworks target general computing environments. While some of these frameworks support multi-node scaling through external integration, they typically assume that the state vector fits entirely within physical GPU VRAM or host DRAM~\cite{cirq, qiskit_aer, qsim, bayraktar2023cuquantum, pennylane}. Several works target HPC systems and use multi-node execution to scale memory capacity, but they rely heavily on high-bandwidth interconnects such as NVLink and InfiniBand~\cite{xu2024atlas, 10.1145/3771577, zhang2021hyquas}. Other works leverage tiered memory and out-of-core data compression for datacenter environments, but target server-grade GPUs~\cite{park2022snuqs, zhang2025bmqsim}. Specifically, while \textit{BMQSim}~\cite{zhang2025bmqsim} scales capacity through tiered memory pipelines, it does not provide the UMA-aware optimizations required to eliminate redundant memory copies on shared-memory edge SoCs. Thus, these existing approaches do not consider the hardware characteristics of edge systems, including restricted NVMe bandwidth and the physical constraints of Unified Memory Architecture (UMA), where traditional host-driven I/O pipelines saturate the memory controller.


In this paper, we present \EdgeQuantum, a scalable full state-vector simulation framework for resource-constrained edge devices. While existing out-of-core simulators address memory limitations through OS-native swapping or server-grade abstractions, they do not coordinate the end-to-end I/O path for UMA environments. \EdgeQuantum integrates unified DRAM and commodity NVMe storage into a 2-tier memory architecture that decouples logical state-vector management from physical residency. Specifically, \EdgeQuantum introduces 1) an out-of-core state vector chunking mechanism that extends logical capacity to NVMe storage beyond hardware limits, 2) an edge-aware asynchronous Unified Virtual Memory (UVM) pipeline that serializes storage access and overlaps flash operations with GPU computation, and 3) an on-the-fly mantissa-truncated LZ4 ping-pong compression strategy that circumvents in-place update limitations through dual-file alternation, achieving up to 255$\times$ per-file compression. Our evaluation on an 8GB Jetson Orin Nano demonstrates that \EdgeQuantum supports up to a 39-qubit simulation (4TB footprint) and achieves a 512$\times$ logical capacity expansion beyond physical DRAM. We open-source \EdgeQuantum at \url{https://github.com/sunggonkim/IoTJ-EdgeQuantum}





%Specifically, \EdgeQuantum introduces 1) state vector partitioning across shared DRAM and NVMe to enable large-scale simulations beyond hardware limits, 2) an edge-aware asynchronous Unified Virtual Memory (UVM) pipeline that serializes storage access to avoid I/O contention and overlaps flash operations with GPU computation, and 3) an on-the-fly LZ4 ping-pong compression strategy that bypasses the impossibility of in-place updates for variable-sized data by alternating backing files, thereby expanding effective storage capacity by $242.7\times$ while masking NVMe latency.